% ============================================
% Supplementary Note 1: Null Models and Statistical Controls
% ============================================

\subsection*{Supplementary Note 1: Null Models and Statistical Controls}
\label{sec:suppl:note1}

A potential concern is that one-sided community-level selection could arise from skewed abundance distributions rather than ecological interactions. To address this, we compared experimental one-sided selection values against two null models.

In Null Model 1 (abundance-weighted random selection), species survival probability is proportional to abundance in the combined parental pool, testing whether high-abundance species simply survive regardless of parental origin. In Null Model 2 (shuffled abundance), abundances are randomly permuted among species within each parent before neutral mixing, testing whether the overall skewness structure alone drives one-sided selection.

We generated 500 null offspring communities per model and calculated the one-sided selection metric. Experimental data showed mean one-sided selection of 0.698 ($n = 265$). Both null models produced significantly lower values: abundance-weighted null mean = 0.476 ($p < 10^{-18}$), shuffled null mean = 0.557 ($p < 10^{-11}$; Mann-Whitney U tests). These results indicate that abundance skewness alone cannot explain the observed patterns; ecological interactions beyond simple abundance-weighted survival determine which parental community dominates (Supplementary Fig. 5).

